\documentclass[10pt,a4paper]{amsart}
\usepackage[margin=19mm]{geometry}
\usepackage{amsmath,amssymb,mathtools}
\usepackage[scr=rsfs,cal=euler]{mathalfa}
\usepackage{xfrac}
\usepackage[unicode,hidelinks,bookmarks=false]{hyperref}
\newcommand{\ie}{i.e.\ }
\newcommand{\cf}{cf.\ }
\newcommand{\resp}{respectively\ }
\newcommand{\Subsection}{\subsection}
\providecommand{\nobreakdash}{\nobreak\hbox{-}}
\setlength{\emergencystretch}{2em}
\multlinegap0pt
\usepackage{bm}
\usepackage{bbm}
\usepackage{dsfont}
\usepackage{xspace}
\usepackage{cancel}
\usepackage{mathtools}
\usepackage{amsthm}
\makeatletter
\@ifundefined{lemma}{\newtheorem{lemma}{Lemma}}{}
\makeatother
\title[Stability of admissible solutions for coexisting phase trans]{Stability of admissible solutions for coexisting phase transitions for one-dimensional compressible van der Waals fluids}
\author{Yazhou Chen, Qiaolin He, Dongjuan Niu, Yi Peng, Xiaoding Shi}
\date{Source: arXiv:2608.20060v1 (CC BY 4.0)}
\begin{document}
\maketitle

\noindent\emph{Source and licence.} Stability of admissible solutions for coexisting phase transitions for one-dimensional compressible van der Waals fluids. Version arXiv:2608.20060v1 (CC BY 4.0), distributed under the Creative Commons Attribution 4.0 International licence, as recorded in the arXiv licence metadata for this exact version and shown on the abstract page https://arxiv.org/abs/2608.20060v1. Full rights notice: licenses/arxiv-2608.20060-CC-BY-4.0.txt.\par\smallskip

\noindent\textit{Editorial scope.} Counted unit: the Lemma whose source label is \texttt{thm-exsit2} in the article (source0.tex lines 555--570, source label \texttt{thm-exsit2}), delivered together with the article's own complete original proof of it (source0.tex lines 572--898, one \texttt{proof} environment of 327 lines). No mathematics is written, repaired or re-derived: statement and proof are the article's own text line for line. Required context retained uncounted: eq:vvv (lines 201--206); eq:scheme (lines 516--521). Every definition, display and label of those lines is retained unchanged. Omitted from this package and not redistributed: 1--200, 207--515, 522--554, 571--571, 899--1548. Nothing inside a retained excerpt is omitted or altered: every retained body is delivered exactly as the article's lines are written, so no omission receipt of any kind accompanies this package. Citations: the retained excerpts carry no citation command at all, so no bibliography entry of the article is transcribed and no reference is redistributed. Reference adaptation: none was needed and none was made; every reference inside the delivered unit resolves inside it, so no unresolved reference or citation is produced. Printed slips kept literal: no printed word, symbol, hypothesis or number of the counted statement or of its proof was altered. Layout and class: the article's own class is \texttt{article} with packages including amsmath, amsfonts, amssymb, amsthm, mathrsfs, mathtools, titlesec, graphicx, subfigure, soul, tocbibind, cancel, fancyhdr, microtype; the wrapper is the lane's pinned \texttt{amsart} editorial wrapper at 10pt on A4, which restates the theorem-like environments the retained text needs and adds the article's own macro definitions that the retained text needs (none), with extra wrapper packages (bm, bbm, dsfont, xspace, cancel, mathtools). The delivered numbering is the unit's own. Assets: no retained span contains a figure environment, a graphics-inclusion command, a PDF-page inclusion, a table or a TikZ picture; no image file is copied into this package, the package declares no asset dependency and no image file is redistributed with it. Licence: version arXiv:2608.20060v1 of this article is distributed under the Creative Commons Attribution 4.0 International licence, as recorded in the arXiv licence metadata for this exact version and shown on the abstract page https://arxiv.org/abs/2608.20060v1. A full rights notice is delivered at licenses/arxiv-2608.20060-CC-BY-4.0.txt. Characters: every retained excerpt is pure ASCII, so the delivered engine, XeLaTeX with its default Latin Modern text font, reports no missing character.
\begin{equation}\label{eq:vvv}
\begin{aligned}
\displaystyle & -C^{-1}\leq p'(v)\leq C,\quad %\\& 
-C^{-1}\leq vp'(v)\leq C,
\end{aligned}	
\end{equation}

\begin{equation}\label{eq:scheme}\left\{
\begin{aligned} 
\displaystyle  &\dot{v_j}=\delta u_j,&&\quad j=\pm \frac12,\pm \frac32,\dots,\pm(N-\frac12),\\
&\dot{u_k}=-\delta p_k+\delta\left(\frac{\epsilon \delta u}{v} \right)_k, &&\quad k=0,\pm 1,\dots,\pm N,
\end{aligned}\right.
\end{equation}

\begin{lemma}\label{thm-exsit2}
Let \(\theta, \underline{v}, \hat{v}\) and \(M_{0}\) be given positive numbers. Then there are constants \(C_1, C_2\), \(t_{0}\), and \(h_0\), depending only on $\theta, \underline{v},\hat v$ and $M_{0}$, such that: given \(h\leq h_{0}\) and initial data \(v_{j}(0), u_{k}(0)\) with values \(b+2\underline{v}\leq v_{j}( 0)\leq \hat v-\underline{v}\)
and with \(\mathscr{C} _{0}\leq M_{0}\), the scheme \eqref{eq:scheme} is solvable up to time \(t_0\) and satisfies
\begin{align}
\displaystyle  &b+\underline{v}\leq v_{j}\leq \hat v, \label{local-01} \\
&\mathscr{A}(t_0)+\mathscr{B}(t_0)\leq C_1\mathscr{C}_0, \label{local-02} %\\&\mathscr{A}(t_0)+\mathscr{D}(t_0)\leq C_2(\mathscr{C}_0+\mathscr{C}_1),\label{local-03}
\end{align}
and the regularity conditions
\begin{align}
&\langle v_{j}\rangle_{I_i\times[0,t_{0}]}^{\frac{1}{2},\frac{1}{2}-\frac{\theta}{4}}\leq C\mathscr{C}_{0},\quad i=1,2,3, \label{local-04}\\
%&\langle u_{k}\rangle_{\{0\leq t\leq t_{0}\}}^{\frac{1}{2},\frac{1}{4}}\leq C(\mathscr{C}_0+\mathscr{C}_1)^{1/2},\label{local-05}\\
&\langle u_k\rangle_{\{\tau\leq t\leq t_0\}}^{\frac{1}{2},\frac{1}{4}}\leq C\mathscr{C}_0\tau^{-\frac{1}{4}-\frac{\theta}{2}},\label{local-06}\\
&\left|\left[\ln v_{\pm N_* }(t)\right]-\exp\left(\epsilon^{-1}\int_{0}^{t}\alpha_{\pm}(s)ds\right)\left[\ln v_{\pm N_* }(0)\right]\right|\leq Ch^{1/2}\mathscr{C}_0^{1/2},\label{local-07}
\end{align}
where \(\alpha_{\pm}(s)\) is given by \eqref{eq:l3}.
\end{lemma}

\begin{proof}
Observe that if \(\mathscr{C}_{0}<\infty\) and $v_{j}(0)\in [b+2\underline{v},\hat{v}-\underline{v}] $, then \(v_j(0)\) and \(u_k(0)\) are bounded uniformly in \(j\) and \(k\) (with bounds depending on \(h\) ). Thus the ordinary  differential equations \eqref{eq:scheme} together with these initial values constitute a well-posed initial value problem in \(L^{\infty}\cap L^{2}\) and so have a local solution defined up to some \(t_0>0\) which may depend on $h$. Our goal is to show that \(t_{0}\) is in fact independent of \(h\). We therefore assume that $v_{j}(t)\in [b+\underline{v},\hat{v}]$ for all $j$ and all $t\in [0,t_0]$ for some positive $t_0$, and proceed to obtain bounds for $\mathscr{A}, \mathscr{B}$ and $\mathscr{D}$. 
Throughout this proof $C$ denotes a generic positive constant as described above. 

We first give the  $L^2$ estimate of the local solution to \eqref{eq:scheme}. We multiply \eqref{eq:scheme}$_{1,2}$ by \((v_j-\bar v)h\) and \((u_k-\bar u)h\), respectively, and sum to obtain
\begin{align*}
\displaystyle  &\frac12\frac{d}{dt}\Big(\sum_{j} (v_j-\bar v)^{2}h+\sum_{k}(u_k-\bar u)^2h\Big)\\&=\sum_{j}\delta u_j(v_j-\bar v)h+\sum_{k}\Big(-\delta(p-\bar p)_k(u_k-\bar u)h+\delta\big(\frac{\epsilon\delta u}{v}\big)_k(u_k-\bar u)h\Big)\\
&=\sum_{j}\delta u_j(v_j-\bar v)h+\sum_{j}\Big((p_j-\bar p)\delta u_jh-\frac{\epsilon\delta u_j^2}{v_j}h\Big),
\end{align*}
where $\bar{p}=p({\bar v})$. Integrating the above equation and using the assumption $b+\underline{v}\leq v_j\leq \hat{v}$, we have
\begin{equation}\label{eq:b1}
\begin{aligned}
\displaystyle  &\Big(\frac12\sum_{j} (v_j-\bar v)(t)^{2}h+\frac12\sum_{k}(u_k-\bar u)(t)^2h\Big)\Big|^t_0+\int_0^t \sum_{j}\delta u_j(s)^2hds\\
&\leq C\int_0^t\sum_{j}\Big(\delta u_j(v_j-\bar v)h+(p_j-\bar p)\delta u_jh\Big)ds.\\
\end{aligned}  
\end{equation}
Noting that \(p_j-\bar p=p'(v_*)(v_j-\bar v)\) and \eqref{eq:vvv}, we can bound the right side of \eqref{eq:b1} by
\[\begin{aligned}
&C\int_0^t\sum_{j}\Big(\delta u_j(v_j-\bar v)h+(p_j-\bar p)\delta u_jh\Big)ds\\& \leq C\int_{0}^{t}\sum_{j}|v_{j}-\bar v||\delta u_j|hds \\&\leq C\int_{0}^{t}\sum_{j}(v_{j}-\bar v)(s)^2hds+\frac12\int_{0}^{t}\sum_{j}\delta u_j(s)^2hds\\&\leq C\mathscr{A}t+\frac12\int_{0}^{t}\sum_{j}\delta u_j(s)^2hds.
\end{aligned}\]
Substituting the above formula into \eqref{eq:b1} yields
\begin{equation}\label{eq:b0}
  \sum_{j}(v_{j}-\bar v)(t)^{2}h+\sum_{k}(u_{k}-\bar u)(t)^{2}h+\int_{0}^{t}\sum\limits_{j}\delta u_{j}(s)^{2}hds\leq C({\mathscr{C}}_{0}+{\mathscr{A}}t).
\end{equation} 

Next, we derive an upper bound for \(\sum\limits_{k^{\prime}}\delta v_{k}^{2}h\).
Letting \(L_j=L(v_j)=\ln v_j\), we find from \eqref{eq:scheme}$_1$ that
\begin{equation}\label{eq:vx1}
  \dot{L}_{j}=\frac{\dot v_{j}}{v_{j}}=\frac{\delta u_{j}}{v_{j}},
\end{equation}
thus we have from \eqref{eq:scheme}$_2$ and \eqref{eq:vx1} that
\begin{equation}\label{eq:vx2}
  \epsilon\delta\dot{L}_{k}=\delta\Big(\frac{\epsilon\delta u}{v}\Big)_{k}=\delta p_k+\dot{u}_k.
\end{equation}
It is easy to obtain from the assumption $b+\underline{v}\leq v_j\leq \hat{v}$ and \eqref{eq:vvv} that
\begin{equation}\label{eq:lpv}
	\delta L_k\sim\delta p_{k}\sim\delta v_{k}.
\end{equation}
Multiplying \eqref{eq:vx2} by \(\delta L_{k}\), summing and integrating, we obtain
\begin{equation*}
\begin{aligned}\frac\epsilon2\sum_{k^{\prime}}\delta L_{k}^{2}h\Big|_{0}^{t}&=\int_{0}^{t}\sum_{k^{'}}\delta L_{k}\delta p_{k}hds+\sum_{k^{'}}(u_k-\bar u)\delta L_{k}h\Big|_{0}^{t}%\\&
-\frac1\epsilon\int_{0}^{t}\sum_{k^{'}} (u_k-\bar u)(\dot{u}_{k}+\delta p_{k})hds\\&
=-\frac1{2\epsilon}\sum_{k^{\prime}} (u_k-\bar u)^{2}h\Big|_{0}^{t}+\int_{0}^{t}\sum_{k^{'}}\delta L_{k}\delta p_{k}hds+\sum_{k^{'}}(u_k-\bar u)\delta L_{k}h\Big|_{0}^{t}\\&\quad
-\frac1\epsilon\int_{0}^{t}\sum_{k^{'}} (u_k-\bar u)\delta p_{k}hds\\
&\leq C\mathscr{C}_{0}+C\sum_{k^{\prime}} (u_k-\bar u)^{2}h+\frac\epsilon4\sum_{k^{\prime}}\!\delta L_{k}^{2}h+\int_{0}^{t}\!\sum_{k^{'}}\!\delta v_{k}^{2}hds+\int_{0}^{t}\!\sum_{k^{'}}(u_{k}\!-\!\bar u)^{2}hds\\
&\leq C\mathscr{C}_{0}+C\sum_{k^{\prime}} (u_k-\bar u)^{2}h+\frac\epsilon4\sum_{k^{\prime}}\!\delta L_{k}^{2}h+\mathscr{A}t,\end{aligned}  
\end{equation*}
which follows from the definition of $\mathscr{A}$, \eqref{eq:vx2} and \eqref{eq:lpv}. Elementary estimates based upon \eqref{eq:b0} then show that
\begin{equation}\label{eq:vx0}
\begin{aligned}\sum_{k'}\delta v_{k}(t)^{2}h\leq C(\mathscr{C}_{0}+\mathscr{A}t).\end{aligned}  
\end{equation}
We obtain from \eqref{eq:b0} and \eqref{eq:vx0} that
 \begin{equation}\label{eq:est-a0}
  \mathscr{A}(t)\leq C(\mathscr{C}_0+\mathscr{A}t).
\end{equation}
Moreover, we find from \eqref{eq:vx2} that at \(x_{\pm N_* }\),
\begin{equation}\label{eq:l1}
  \epsilon[\dot{L}_{\pm N_* }]=\Big[\big(\frac{\epsilon\delta u}{v}\big)_{\pm N_* }\Big]=[p_{\pm N_* }]+h\dot{u}_{\pm N_* },
\end{equation}
where \([w_{\pm N_* }]\triangleq w_{\pm N_* +\frac12}-w_{\pm N_* -\frac12}\). We obtain
\begin{equation}\label{eq:l2}
 [p_{\pm N_* }]=\alpha_\pm[L_{\pm N_* }],
\end{equation}
in which 
\begin{equation}\label{eq:l3}
 \alpha_{\pm}=\frac{p(v_{\pm N_* +\frac12})-p(v_{\pm N_* -\frac12})}{L(v_{\pm N_* +\frac12})-L(v_{\pm N_* -\frac12})}=\xi_\pm p'(\xi_\pm),
\end{equation}
with $\xi_\pm$ are between $v_{\pm N_* -\frac12}$ and $v_{\pm N_* +\frac12}$. 
The bounds in \eqref{eq:vvv} show that $-C^{-1}\leq\alpha_\pm\leq C$. Equation \eqref{eq:l1} thus becomes
\begin{equation}\label{eq:l4}
  \epsilon[\dot{L}_{\pm N_* }]=\alpha_\pm[L_{\pm N_* }]+h\dot{u}_{\pm N_* },
\end{equation}
whose solution is
\[\begin{aligned}
[L_{\pm N_* }(t)]& =\exp\Big(\epsilon^{-1}\int_{0}^{t}\alpha_{\pm}(s)ds\Big)\Big([L_{\pm N_* }(0)]-\frac{h}{\epsilon}(u_{\pm N_* }(0)-\bar u)\Big)+\frac{h}{\epsilon}(u_{\pm N_* }(t) -\bar u) \\
&+\frac{h}{\epsilon^{2}}\int_{0}^{t}\alpha_{\pm}(s)\exp\Big(\epsilon^{-1}\int_{s}^{t}\alpha_{\pm}(\tau)d\tau\Big)(u_{\pm N_* }(s) -\bar u)ds.
\end{aligned}\]
The bound \(|u_{\pm N_* }-\bar u|\leq h^{-1/2}\mathscr{A}^{1/2}\) and the assumption $0\leq t\leq t_0\leq 1$ show that
\begin{equation}\label{eq:lnv-jump0}\Big|[L_{\pm N_* }(t)]-\exp\Big(e^{-1}\int_{0}^{t}\alpha_{i}(s)ds\Big)[L_{\pm N_* }(0)]\Big|\leq Ch^{\frac12}\mathscr{A}^{\frac12},\end{equation}
so that by the above bound for \(\alpha_\pm\), it holds
\begin{equation}\label{eq:v-jump0}\left|[v_{\pm N_* }(t)]\right|\leq C(\left|[v_{\pm N_* }(0)]\right|+h^{\frac12}\mathscr{A}^{\frac12}).\end{equation}
We also have the estimate
\begin{equation}\label{eq220}
|[p_{\pm N_* }(t)]|\leq C(\tilde{\mathscr{C}}^{1/2}+h^{\frac12}\mathscr{A}^{\frac12}),
\end{equation}
which follows from \eqref{eq:l2} and \eqref{eq:v-jump0}. 



We now proceed to derive a bound for $\mathscr{B}$. 
We multiply the second equation in \eqref{eq:scheme} by \(\delta\psi_{k}(s)h\)  where \(\{\psi_{j}(s)\}_{0\leq s\leq t}\) is a test sequence to be specified below, sum the product over \(j\) and integrate theresult to obtain
\[-\sum_{j}\psi_{j}\delta u_{j}h\Big|_{0}^t=-\int_{0}^{t}\sum_{j}\delta u_{j}\left(\dot{\psi}_{j}+\frac{\epsilon\delta(\delta\psi)_{j}}{v_{j}}\right)hds-\int_{0}^{t}\sum_{k}\delta\psi_{k}\delta p_{k}hds.\]
We now choose \(\psi_{j}\) to satisfy the initial value problem
\begin{equation}\label{back-para}
  \left\{\begin{aligned}
  \displaystyle  &\dot{\psi}_{j}+\frac{\epsilon\delta(\delta\psi)_{j}}{v_{j}}=0,\quad0\leq s\leq t,\\&\psi_{j}(t)=\delta u_{j}(t)\Big/\Big(\sum_{j}\delta u_{j}(t)^{2}h\Big)^{1/2}.
  \end{aligned}  \right.
\end{equation}
We then have that
\begin{equation}\label{eq228}
 \begin{aligned}
\Big(\sum_{j}\delta u_{j}(t)^{2}h\Big)^{1/2}\leq & \Big|\sum_{j}\psi_{j}(0)\delta u_{j}(0)h\Big|+\int_{0}^{t}\sum_{k^{\prime}}|\delta\varphi_{k}\delta p_{k}|hds \\&+\int_{0}^{t}\|\delta\psi(s)\|_{\infty}(|[p_{N_* }(s)]|+|[p_{- N_* }(s)]|)ds\\\triangleq &\sum_{i=1}^3I_i.
\end{aligned} 
\end{equation}
Defining the quantity $\mathscr{E}$ by
\[\begin{aligned}\label{eq229}
\mathscr{E}=& \sup_{0\leq s\leq t}\Big(\sum_{j}\psi_{j}(s)^{2}h+(t-s)\sum_{k}\delta\psi_{k}(s)^{2}h\Big)  \\
&+\int_{0}^{t}\Big(\sum\limits_{k}\delta\psi_{k}(s)^{2}h+(t-s)\sum_{j}(\dot{\psi}_{j}(s)^{2}+\delta(\delta\psi)_{j}(s)^{2})h\Big)ds.
\end{aligned}\]
We proceed to estimate the right-hand side of \eqref{eq228} in terms of $\mathscr{E}$. First,
\[\|\psi(0)\|_{\infty}\leq \Big(\sum_{j}\psi_{j}(0)^{2}h\Big)^{1/4}\Big(\sum_{k}\delta\psi_{k}(0)^{2}h\Big)^{1/4}\leq Ct^{-1/4}\mathscr{E}^{1/2},\]
so that
\[\begin{aligned}
I_1=\Big|\sum\psi_{j}(0)\delta u_{j}(0)h \Big|&\leq Ct^{-1/4}\mathscr{E}^{1/2}\sum_{j}\mid u_{j+\frac{1}{2}}(0)-u_{j-\frac{1}{2}}(0)\mid \\
&\leq Ct^{-1/4}\mathscr{E}^{1/2}\mathscr{C}_{0}^{1/2}. 
\end{aligned}\]
Next, from \eqref{eq:lpv} and \eqref{eq229},
\[\begin{aligned}
I_2=\int_{0}^{t}\sum_{k^{\prime}}|\delta\varphi_{k}|~|\delta p_{k}| hds &\leq C\Big(\int_{0}^{t}\sum\limits_{k^{\prime}}\delta v_{k}^{2}hds\Big)^{1/2}\mathscr{E}^{1/2} \leq C\mathscr{A}^{1/2}\mathscr{E}^{1/2}.
\end{aligned}\]
Finally, using \eqref{eq229} and \eqref{eq220}, we have
\[\begin{aligned}
I_3&\leq\int_{0}^{t}\Big(\sum_{k}\delta\psi_{k}^{2}h\Big)^{1/4}\Big(\sum_{j}\delta(\delta\psi)_{j}^{2}h\Big)^{1/4}(\tilde{\mathscr{C}}^{1/2}+h^{1/2}{\mathscr{A}}^{1/2})ds \\
&\leq C\mathscr{E}^{\frac14}\Big(\int_{0}^{t}(t-s)^{-\frac23}ds\Big)^{\frac34}\Big(\int_{0}^{t}\sum_{j}(t-s)\delta(\delta\psi)_{j}^{2}hds\Big)^{\frac14}(\tilde{\mathscr{C}}^{\frac12}+h^{\frac12}{\mathscr{A}}^{\frac12})\\
&\leq C\mathscr{E}^{1/2}t^{\frac14}(\tilde{\mathscr{C}}^{\frac12}+h^{\frac12}{\mathscr{A}}^{\frac12}).
\end{aligned}\]
Substituting the above estimates into \eqref{eq228} yields
\begin{equation}\label{eq230}
\begin{aligned}
  \sup\limits_{0\leq t\leq t_{0}}t^{1/2}\sum\limits_{j}\delta u_{j}(t)^{2}h&\leq C{\mathscr{E}}(\mathscr{C}_{0}+\mathscr{A}t^{1/2}+\tilde{\mathscr{C}}t+ h\mathscr{A}t)\\
  &\leq C{\mathscr{E}}(\mathscr{C}_{0}+\mathscr{A}t_0^{1/2}+ h\mathscr{A}),
\end{aligned}
\end{equation}
by the assumption $\tilde{\mathscr{C}}t\leq \mathscr{C}_{0}$ for small time $t\leq t_0$.

We now derive  the estimate of $\mathscr{E}$ in terms of $\mathscr{A}$.
First, we multiply \eqref{back-para} by \(v_j\psi_jh\), sum and integrate to obtain
$$\sum_{j}\psi_{j}(s)^{2}h+\int_{s}^{t}\sum\limits_{k}\delta\psi_{k}(\tau)^{2}hd\tau\leq C\Big(1+\int_{s}^{t}\sum_{j}\psi_{j}^{2}|\delta u_{j}|hd\tau\Big).$$
The last term can be bounded by
\begin{equation*}
\begin{aligned}
\int_{s}^{t}\sum_{j}\psi_{j}^{2}|\delta u_{j}|hd\tau
&\leq\int_{s}^{t}\Big(\sum_{j}\psi_{j}(\tau)^{2}h\Big)^{1/2}\Big(\sum_{j}\delta u_{j}(\tau)^{2}h\Big)^{1/2}\|\psi_j\|_\infty d\tau\\
&\leq C{\mathscr{E}}\int_{s}^{t}\Big(\sum_{j}\delta u_{j}(\tau)^{2}h\Big)^{1/2}(t-\tau)^{-\frac14} d\tau\\
&\leq C{\mathscr{E}}{\mathscr{A}}^{1/2}t^{1/4},
\end{aligned}
\end{equation*}
thus
\begin{equation}\label{eq233}
\begin{aligned}
\sum_{j}\psi_{j}(s)^{2}h+\int_{s}^{t}\sum\limits_{k}\delta\psi_{k}(\tau)^{2}hd\tau
&\leq C(1+{\mathscr{E}}{\mathscr{A}}^{1/2}t^{1/4}).
\end{aligned}
\end{equation}
Multiplying \eqref{back-para} by \(\left(t-s\right)v_{j}\dot{\psi}_{j}h\), summing
and integrating, we obtain in a similar way that
\begin{equation}\label{eq234}
(t-s)\sum\limits_{k}\delta\psi_{k}(s)^{2}h+\int_{s}^{t}\sum\limits_{j}(t-\tau)\dot{\psi}_{j}(\tau)^{2}hd\tau\leq C\int_{s}^{t}\sum\limits_{k}\delta\psi_{k}(\tau)^{2}hd\tau.\end{equation}
Combining appropriate multiples of \eqref{eq233}, \eqref{eq234} and \eqref{back-para}, and taking the
supremum over \(s\in[0,t]\), we obtain
\begin{equation}\label{eq231}
\mathscr{E}\leq C(1+\mathscr{E}\mathscr{A}^{1/2}t^{1/4}).  
\end{equation}
It holds that \(\mathscr{A}t^{1/2}\leq1/(4C^{2})\) for small time $t\leq t_0$, which shows \(\mathscr{E}\leq C\). Substituting this into \eqref{eq230}, we obtain 
\begin{equation}\label{eq:est-ux}
\sup_{0\leq t\leq t_{0}}t^{1/2}\sum_{j}\delta u_{j}(t)^{2}h\leq C(\mathscr{C}_{0}+\mathscr{A}t_0^{1/2}+ h\mathscr{A}).\end{equation}

Next, we multiply the second equation in \eqref{eq:scheme} by \(t^{\beta}\dot{u}_{k}(t)h\), sum and integrate to get
\begin{equation}\label{eq:t-beta}\begin{aligned}
		&\frac{s^{\beta}}{2}\sum_{j}\frac{\delta u_{j}(s)^{2}h}{v_j}\Big|_0^t+\int_{0}^{t}\sum_{k}s^{\beta}\dot{u}_{k}(s)^{2}hds \\&= s^{\beta}\sum_{j}(p_j-\bar p)\delta u_j h\Big|_0^t+\int_0^t\frac{\beta s^{\beta-1}}{2}\sum_{j}\frac{\delta u_{j}(s)^{2}h}{v_j}ds%&\\
		-\int_0^t\frac{s^{\beta}}{2}\sum_{j}\frac{\delta u_{j}(s)^{3}h}{v_j^2}ds \\&
		\quad-\int_{0}^{t}s^{\beta}\sum_{j}\dot{p_j}\delta u_j hds-\int_{0}^{t}\beta s^{\beta-1}\sum_{j}(p_j-\bar p)\delta u_j hds
\end{aligned}\end{equation}
Let $\beta=\frac{1}{2}+\theta$ with \(0<\theta<1\), and after some elementary manipulations, we obtain
\begin{equation}\label{eq235}\begin{aligned}
&t^{\frac{1}{2}+\theta}\sum_{j}\delta u_{j}(t)^{2}h+\int_{0}^{t}\sum_{k}s^{\frac{1}{2}+\theta}\dot{u}_{k}(s)^{2}hds \\
&\leq t^{\frac{1}{2}+\theta}\sum_{j}|p_j-\bar p|\,|\delta u_j| h+C\int_0^ts^{-\frac{1}{2}+\theta}\sum_{j}\delta u_{j}(s)^{2}hds\\
&\quad		+C\int_0^ts^{\frac{1}{2}+\theta}\sum_{j}|\delta u_{j}(s)|^{3}hds +\int_{0}^{t}s^{\frac{1}{2}+\theta}\sum_{j}|\dot{p_j}|\,|\delta u_j| hds\\&\quad+C\int_{0}^{t}s^{-\frac{1}{2}+\theta}\sum_{j}|p_j-\bar p|\,|\delta u_j| h \\&\triangleq \sum_{i=1}^{5}J_i.
\end{aligned}\end{equation}
We apply straightforward estimates to each term on the right-hand side of \eqref{eq235}. First, By \eqref{eq:vvv},\eqref{eq:lpv} and the definition of $\mathscr{A},\mathscr{B}$, we obtain
\begin{align}
	&J_1\leq C t^{\frac{1}{2}+\theta}\sum_{j}|v_j-\bar v|\,|\delta u_j| h\leq C\mathscr{A}^{1/2}\mathscr{B}^{1/2}t^{\frac{1}{4}+\theta},\label{eq:ux-j1}\\
	&J_2\leq C\mathscr{B}\int_0^ts^{-1+\theta}ds\leq C\mathscr{B}t^\theta,\\
	&J_4\leq C\int_{0}^{t}s^{\frac{1}{2}+\theta}\sum_{j}|\dot{v_j}|\,|\delta u_j| hds\leq C\mathscr{B}\int_0^ts^{\theta}ds\leq C\mathscr{B}t^{1+\theta},\\
	&J_5\leq C\int_0^ts^{-\frac{1}{2}+\theta}\sum_{j}|v_j-\bar v|\,|\delta u_j| hds%\leq C\mathscr{A}^{1/2}\mathscr{B}^{1/2}\int_0^ts^{-\frac{3}{4}+\theta}ds
	\leq C\mathscr{A}^{1/2}\mathscr{B}^{1/2}t^{\frac{1}{4}+\theta}.\label{eq:ux-j5}
\end{align}
Similarly, we have
\begin{equation}\label{eq:ux-j3}
\begin{aligned}
J_3&\leq s^{\frac{1}{2}+\theta}\int\sum |\delta u_{j}|^{3}hds  \\
&\leq C\int_{0}^{t}s^{\theta}\sum_{j}\delta u_{j}^{2}hds+C\int_{0}^{t}s^{1+\theta}\sum_{j}\delta u_{j}^{4}hds \\
&\leq C\mathscr{B}t^{\frac{1}{2}+\theta}+C\mathscr{B}\int_{0}^{t}s^{\frac12+\theta}\|\delta u\|_{\infty}^2(s)ds.
\end{aligned}	
\end{equation}
Hence we need to obtain the bound for \(\|\delta u\|_{\infty}\). A discrete Sobolev-type inequality shows that
\[\begin{aligned}
\left({\frac{\delta u}{v}}\right)_J^{2}-\left({\frac{\delta u}{v}}\right)_j^{2}& =\sum_{j<k<J}\Big(\big(\frac{\delta u}{v}\big)_{k+\frac{1}{2}}^{2}-\big(\frac{\delta u}{v}\big)_{k-\frac{1}{2}}^{2}\Big) \\
&\leq C\Big(\sum_{j}\delta u_{j}^{2}h\Big)^{1/2}\Big(\sum_{k^{\prime}}\delta\big(\frac{\delta u}{v}\big)_{k}^{2}h\Big)^{1/2} \\
&\quad+C\Big(\eta\left\|\delta u\right\|_{\infty}^{2}+\eta^{-1}\big(\big|\big[\frac{\delta u}{v}\big]_{\pm N_* }\big|\big)^{2}\Big).
\end{aligned}\]
Summing the above inequality over $j$ from $-N+1/2$ to $N-1/2$, we obtain
\[\begin{aligned}
\left({\frac{\delta u}{v}}\right)_J^{2}&\leq C\sum_{j}\delta u_{j}^{2}h+
 C\Big(\sum_{j}\delta u_{j}^{2}h\Big)^{1/2}\Big(\sum_{k^{\prime}}\delta\big(\frac{\delta u}{v}\big)_{k}^{2}h\Big)^{1/2} \\
&\quad+C\Big(\eta\left\|\delta u\right\|_{\infty}^{2}+\eta^{-1}\big(\big|\big[\frac{\delta u}{v}\big]_{N_* }\big|+\big|\big[\frac{\delta u}{v}\big]_{-N_* }\big|\big)^{2}\Big).
\end{aligned}\]
Taking $\eta$ small and applying \eqref{eq:l1} and \eqref{eq220}, we thus obtain
\begin{equation}\label{eq:ux-infty}
\begin{aligned}
\|\delta u(t)\|_{\infty}^{2}\leq Ct^{-\frac12}\mathscr{B}+Ct^{-\frac14}\mathscr{B}^{1/2}\Big(\sum_{k^{\prime}}\delta\big(\frac{\delta u}{v}\big)_{k}^{2}h\Big)^{1/2}+C\big(\tilde{\mathscr{C}}+ h\mathscr{A}+h^{2}(\dot{u}_{N_* }^{2}+\dot{u}_{-N_* }^{2})\big).
\end{aligned}	
\end{equation}
% Also, it holds
% \begin{equation}\label{eq:ux-infty1}
% \begin{aligned}
% \|\delta u(t)\|_{\infty}^{2}\leq C\mathscr{D}+C\mathscr{D}^{1/2}\Big(\sum_{k^{\prime}}\delta\big(\frac{\delta u}{v}\big)_{k}^{2}h\Big)^{1/2}+C\big(\tilde{\mathscr{C}}+ h\mathscr{A}+h^{2}(\dot{u}_{N_* }^{2}+\dot{u}_{-N_* }^{2})\big).
% \end{aligned}	
% \end{equation}
From above estimates, we have
\begin{equation}\label{eq:ux-j31}
\begin{aligned}
&\int_{0}^{t}s^{\frac12+\theta}\|\delta u\|_{\infty}^2(s)ds\\
&\leq C\int_{0}^{t}s^{\theta}\mathscr{B}ds+C\int_{0}^{t}s^{\frac14+\theta}\mathscr{B}^{1/2}\Big(\sum_{k^{\prime}}\delta\big(\frac{\delta u}{v}\big)_{k}^{2}h\Big)^{1/2}ds\\
&\quad+C\int_{0}^{t}s^{\frac12+\theta}\big(\tilde{\mathscr{C}}+ h\mathscr{A}\big)ds+C\int_{0}^{t}s^{\frac12+\theta}h^{2}(\dot{u}_{N_* }^{2}+\dot{u}_{-N_* }^{2})ds\\
&\leq C\mathscr{B}t^{1+\theta}+C\mathscr{B}t^{\frac12+\frac\theta 2}+C\big(\tilde{\mathscr{C}}+ h\mathscr{A}\big)t^{\frac32+\theta}+C\mathscr{B}h.
\end{aligned}	
\end{equation}
Substituting \eqref{eq:ux-j31} into \eqref{eq:ux-j3} yields
\begin{equation}\label{eq:ux-j32}
\begin{aligned}
J_3&\leq  C\mathscr{B}t^{\frac{1}{2}+\theta}+C\mathscr{B}\Big(\mathscr{B}t^{1+\theta}+\mathscr{B}t^{\frac12+\frac\theta 2}+\big(\tilde{\mathscr{C}}+ h\mathscr{A}\big)t^{\frac32+\theta}+\mathscr{B}h\Big)\\
&\leq  C(\mathscr{B}+\mathscr{B}^2)t^{\theta}+C h(\mathscr{A}^2+\mathscr{B}^2).
\end{aligned}	
\end{equation}
By combining \eqref{eq:ux-j1}-\eqref{eq:ux-j5}, \eqref{eq:ux-j32} and \eqref{eq235}, we obtain
\begin{equation}\label{eq236}\begin{aligned}
&t^{\frac{1}{2}+\theta}\sum_{j}\delta u_{j}(t)^{2}h+\int_{0}^{t}\sum_{k}s^{\frac{1}{2}+\theta}\dot{u}_{k}(s)^{2}hds \\
&\leq C\mathscr{A}^{1/2}\mathscr{B}^{1/2}t^{\frac{1}{4}+\theta}+C\mathscr{B}t^{1+\theta}+C(\mathscr{B}+\mathscr{B}^2)t^{\theta}+C h(\mathscr{A}^2+\mathscr{B}^2)\\
&\leq C(\mathscr{A}+\mathscr{B}+\mathscr{B}^2)t^{\theta}+C h(\mathscr{A}^2+\mathscr{B}^2),
\end{aligned}\end{equation}
which yields
\begin{equation}\label{eq:est-uxadd}\begin{aligned}
&\sup_{0\leq t\leq t_{0}}t^{\frac{1}{2}+\theta}\sum_{j}\delta u_{j}(t)^{2}h+\int_{0}^{t_0}\sum_{k}t^{\frac{1}{2}+\theta}\dot{u}_{k}(t)^{2}hdt\\
&\quad\leq C(\mathscr{A}+\mathscr{B}+\mathscr{B}^2)t_0^{\theta}+C h(\mathscr{A}^2+\mathscr{B}^2).
\end{aligned}\end{equation}
By virtue of the relations \eqref{eq:scheme},  \eqref{eq:vx0} and \eqref{eq:est-uxadd}, we obtain
\begin{equation*}
	\begin{aligned}
\int_{0}^{t_{0}}t^{\frac12+\theta}\sum\limits_{k^{\prime}}\delta\left(\frac{\delta u}{v}\right)_{k}(t)^{2}hdt&\leq \int_{0}^{t_{0}}t^{\frac12+\theta}\sum_{k^{\prime}}\dot{u}_{k}(t)^{2}hdt+\int_{0}^{t_{0}}t^{\frac12+\theta}\sum_{k^{\prime}}\delta p_{k}(t)^{2}hdt \\
&\leq C\mathscr{C}_0+C(\mathscr{A}+\mathscr{B}+\mathscr{B}^2)t_0^{\theta}+C h(\mathscr{A}^2+\mathscr{B}^2),
\end{aligned}
\end{equation*}
which together with \eqref{eq:est-uxadd} yields 
\begin{equation}\label{prop-ab}\mathscr{B}(t)\leq C\big(\mathscr{C}_0+(\mathscr{A}+\mathscr{B}+\mathscr{B}^2)t^{\frac{\theta}{2}}+ h(\mathscr{A}+\mathscr{A}^2+\mathscr{B}^2)\big).\end{equation}


% In a similar way, let $\beta=0$ in \eqref{eq:t-beta}, we obtain
% \begin{equation}\label{eq:ux-t0}\begin{aligned}
% &\sum_{j}\delta u_{j}(t)^{2}h+\int_{0}^{t}\sum_{k}\dot{u}_{k}(s)^{2}hds \\
% &\leq C(\mathscr{C}_0+\mathscr{C}_1)+\sum_{j}|p_j-\bar p|\,|\delta u_j| h		+C\int_0^t\sum_{j}|\delta u_{j}(s)|^{3}hds %\\&		\quad
% +\int_{0}^{t}\sum_{j}|\dot{p_j}|\,|\delta u_j| hds\\
% &\leq C(\mathscr{C}_0+\mathscr{C}_1)+\eta\sum_{j}\delta u_{j}(t)^{2}h+\eta^{-1}\mathscr{A}+C\mathscr{D}t^{\frac12}\Big(\int_0^t\|u_{j}\|_\infty^{2}ds\Big)^{\frac12}+C\mathscr{D}t.
% \end{aligned}\end{equation}
% Taking $\eta$ small and applying \eqref{eq:est-a0} and \eqref{eq:ux-infty1}, we thus obtain
% \begin{align}\label{eq:est-uxadd1}
% &\sup_{0\leq t\leq t_{0}}\sum_{j}\delta u_{j}(t)^{2}h+\int_{0}^{t_{0}}\sum_{k}\dot{u}_{k}(t)^{2}hdt \\&\quad
% \leq C(\mathscr{C}_{0}+\mathscr{C}_{1})+C(\mathscr{A}+\mathscr{D}+\mathscr{D}^2)t_0^{\frac12}+C h(\mathscr{A}+\mathscr{D}). 
% \end{align}
% Similarly, by virtue of the relations \eqref{eq:scheme},  \eqref{eq:vx0} and \eqref{eq:est-uxadd1}, we have
% \begin{equation*}
% 	\begin{aligned}
% \int_{0}^{t_{0}}\sum\limits_{k^{\prime}}\delta\left(\frac{\delta u}{v}\right)_{k}(t)^{2}hdt&\leq \int_{0}^{t_{0}}\sum_{k^{\prime}}\dot{u}_{k}(t)^{2}hdt+\int_{0}^{t_{0}}\sum_{k^{\prime}}\delta p_{k}(t)^{2}hdt \\
% &\leq C(\mathscr{C}_{0}+\mathscr{C}_{1})+C(\mathscr{A}+\mathscr{D}+\mathscr{D}^2)t^{\frac12}+C h(\mathscr{A}+\mathscr{D}). 
% \end{aligned}
% \end{equation*}
% which together with \eqref{eq:est-uxadd1} yields 
% \begin{equation}\label{prop-ad}
%   \mathscr{D}(t)\leq C\big(\mathscr{C}_0+\mathscr{C}_1)+(\mathscr{A}+\mathscr{D}+\mathscr{D}^2)t^{\frac12}+ h(\mathscr{A}+\mathscr{D})\big).
% \end{equation}


Let \(C=C(\theta,\underline{v},\hat v,M_0)\) be the constant in \eqref{eq:est-a0}  and \eqref{prop-ab}. For small time $t_0$, we also choose \(h_{0}\) sufficiently small such that \(h\leq h_{0}\) and \(t\leq t_0\), there exists a constant $C_1=C_1(\theta,\underline{v},\hat v,M_0)$ satisfying
\[\mathscr{A}+\mathscr{B}\leq C_1\mathscr{C}_0.\]
Furthermore, by \eqref{eq:ux-infty}, it holds
\begin{equation}\label{eq:v-infty}
	\begin{aligned}
\left| {{v}_{j}\left( {t}\right)  - {v}_{j}\left( 0\right) }\right|  &\leq  \int_{0}^{{t}}\parallel {\delta u}\left( s\right) {\parallel }_{\infty }{ds}\\
&\leq  C\int_{0}^{t}  {s}^{-\frac{1}{4}}{\mathscr{B}}^{\frac{1}{2}}ds+C\int_{0}^{t}  {{s}^{-\frac{1}{4} - \frac{\theta }{4}}{\mathscr{B}}^{\frac{1}{4}}{\left( {s}^{\frac{1}{2} + \theta }\mathop{\sum }\limits_{k^{\prime}}\delta {\left( \frac{\delta u}{v}\right) }_{k}^{2}h\right) }^{\frac{1}{4}}}ds\\
&\quad +C\int_{0}^{t}\big(\tilde{\mathscr{C}}+ h\mathscr{A}+h^{2}\dot{u}_{\pm N_* }(t)^{2}\big)^{\frac12}ds\\
&\leq  C\Big({\mathscr{B}}^{\frac{1}{2}}{t}^{\frac{3}{4}}+  {\mathscr{B}}^{\frac{1}{2}}{t}^{\frac{1}{2} - \frac{\theta }{4}} + (\tilde{\mathscr{C}}+ h\mathscr{A})^{\frac{1}{2}}t+ {\mathscr{B}}^{\frac{1}{2}}{t}^{\frac{3}{4} - \frac{\theta }{2}} \Big)\\
&\leq  C{\mathscr{C}}_{0}^{\frac{1}{2}}t^{\frac{1}{2} - \frac{\theta }{4}}\leq \underline{v},
	\end{aligned}	
\end{equation}
provided that \(t\leq t_0\) and \(h\leq h_{0}=O(t_0)\). Thus
\begin{equation}\label{eq:v-infty1}
	b+\underline{v}\leq v_j(0)-\underline{v}\leq v_j(t)\leq v_j(0)+\underline{v}\leq \hat{v}.
\end{equation}
The proof of \eqref{local-01}-\eqref{local-02} then proceeds as above.

Next, we derive the bounds \eqref{local-04}-\eqref{local-07} as consequences of the energy estimates \eqref{local-02}. This will show that \eqref{local-04}-\eqref{local-07} hold, provided that \(h\leq h_{0}\) and \(t\leq t_0\), and for as long as the solution has values \(v_{j} \in [b+\underline{v}, \hat{v}]\). 
First, the validity of the assertions \eqref{local-04}-\eqref{local-06} of H\"{o}lder continuity in the \(x\)-variable is evident from the energy estimates \eqref{local-02}, and the definitions \(\mathscr{A},\mathscr{B}\). To prove the regularity of \(\left\{{v}_{j}\right\}\) in \(t\), we compute from \eqref{eq:scheme} and \eqref{eq:ux-infty} that
\begin{equation}\label{eq256}
	\begin{aligned}
\left| {{v}_{j}\left( {t}^{\prime }\right)  - {v}_{j}(t) }\right|  &\leq  \int_{t}^{{t}^{\prime }}\parallel {\delta u}\left( s\right) {\parallel }_{\infty }{ds}
% \\
% &\leq C\int_{t}^{{t}^{\prime}}  {s}^{-\frac{1}{4}}{\mathscr{B}}^{\frac{1}{2}}ds+ C\int_{t}^{{t}^{\prime}}  {{s}^{-\frac{1}{4} - \frac{\theta }{4}}{\mathscr{B}}^{\frac{1}{4}}{\left( {s}^{\frac{1}{2} + \theta }\mathop{\sum }\limits_{k^{\prime}}\delta {\left( \frac{\delta u}{v}\right) }_{k}^{2}h\right) }^{\frac{1}{4}}}ds\\
% &\quad +C\int_{t}^{{t}^{\prime}}\big(\tilde{\mathscr{C}}+ h\mathscr{A}+h^{2}\dot{u}_{\pm N_* }(t)^{2}\big)^{\frac12}ds\\
% &\leq  C\Big( {\mathscr{B}}^{\frac{1}{2}}\big| {\left( {t}^{\prime }\right) }^{\frac{3}{4}} - {t}^{\frac{3}{4}}\big|+ {\mathscr{B}}^{\frac{1}{2}}{\big| {\left( {t}^{\prime }\right) }^{\frac{2 - \theta }{3}} - {t}^{\frac{2 - \theta }{3}}\big| }^{\frac{3}{4}} + (\tilde{\mathscr{C}}+ h\mathscr{A})^{\frac{1}{2}}\big| {{t}^{\prime } - t}\big|\\
% &\quad\quad+ (h\mathscr{D})^{\frac{1}{2}}\big| {({t}^{\prime}) }^{\frac{1}{2}} - {t}^{\frac{1}{2}}\big| \Big)\\&
\leq  C \mathscr{C}_{0}^{\frac{1}{2}}{\big| {t}^{\prime } - t\big| }^{\frac{1}{2} - \frac{\theta }{4}}.
	\end{aligned}	
\end{equation}
We have taken \(h\) and \(t\) small here, depending on \({M}_{0}\).
To prove the H\"{o}lder continuity \eqref{local-06} for $u$, we first observe that, for \(J \in  \mathbb{Z}\) and \(s \geq  \tau\),
\[
\left| {u_k(s)  - \frac{1}{Jh}\mathop{\sum }\limits_{{l = k}}^{{k + J - 1}}{u}_{k}\left( s\right) h}\right|  \leq  \mathop{\sum }\limits_{{j = k + 1/2}}^{{k + J - 3/2}}\left| {\delta {u}_{j}}\right| h \leq  \tau^{\frac14}\mathscr{B}^{1/2}(Jh)^{1/2}.
\]
Thus
\[\begin{aligned}
\displaystyle 
\left| {{u}_{k}\left( {t}^{\prime }\right)  - {u}_{k}(t) }\right| & \leq  \frac{1}{Jh}\int_{t}^{{t}^{\prime }}\mathop{\sum }\limits_{k}^{{k + J - 1}}\left| {{\dot{u}}_{l}\left( s\right) }\right| {hds} + 2\tau^{-\frac14}\mathscr{B}^{1/2}(Jh)^{1/2}\\
&\leq C\tau^{-\frac14-\frac{\theta}{2}}\mathscr{B}^{1/2}(Jh)^{-1/2}|t'-t|^{\frac12}+2\tau^{-\frac14}\mathscr{B}^{1/2}(Jh)^{1/2}\\
& \leq  C\tau^{-\frac14-\frac{\theta}{2}}\mathscr{C}_0^{1/2}{\left| {t}^{\prime } - t\right| }^{\frac14},
\end{aligned}\]
if we take \({Jh} = O( {| {t}^{\prime } - t| }^{1/2})\), this proves \eqref{local-06}. Finally, \eqref{local-07} follows from \eqref{eq:lnv-jump0} and \eqref{local-02}. It completes the proof of Lemma \ref{thm-exsit2}.
\end{proof}

\end{document}
